\documentclass{article}
\usepackage{amsmath, amssymb, amsthm}

\newtheorem{theorem}{Theorem}[section]
\newtheorem{lemma}[theorem]{Lemma}
\newtheorem{proposition}[theorem]{Proposition}
\newtheorem{corollary}[theorem]{Corollary}
\newtheorem{conjecture}[theorem]{Conjecture}
\theoremstyle{definition}
\newtheorem{definition}[theorem]{Definition}

\title{Erdős Problem \#1196: Primitive Sets and Logarithmic Density}
\author{Fermat Test Harness}
\date{}

% ── Context ────────────────────────────────────────────────────────────────
%   Erdős Problem #1196 asks whether the primitive-set sum
%     S(A) = sum_{a in A} 1 / (a log a)
%   is bounded by 1 + o(1) whenever A lies in [x, infty). The bound is trivial
%   for primes (by Mertens), and Erdős proved in 1935 that S(A) is uniformly
%   bounded over all primitive A ⊂ N. Lichtman later pushed the constant
%   down to ~1.399. The full conjecture was recently resolved: for any
%   primitive A ⊂ [x, infty), S(A) ≤ 1 + O(1/log x). Formalized in Lean.
%
%   This file is a graduated test for Fermat's proof pipeline:
%     • Easy   → definitions + trivial observations
%     • Medium → Erdős 1935 boundedness + Mertens on primes
%     • Hard   → the main Erdős-Sárközy-Szemerédi 1+o(1) conjecture itself
% ────────────────────────────────────────────────────────────────────────────

\begin{document}
\maketitle

\section{Primitive Sets}

\begin{definition}[Primitive Set]
\label{def:primitive}
A set $A \subseteq \mathbb{N}_{\geq 2}$ is \emph{primitive} if no element of
$A$ properly divides another; equivalently, for all distinct $a, b \in A$
we have $a \nmid b$.
\end{definition}

\begin{definition}[Primitive-set sum]
\label{def:sum-S}
For a primitive set $A \subseteq \mathbb{N}_{\geq 2}$, define
\[
  S(A) \;:=\; \sum_{a \in A} \frac{1}{a \log a}.
\]
\end{definition}

\begin{lemma}[Set of primes is primitive]
\label{lem:primes-primitive}
The set of prime numbers $\mathbb{P} = \{2, 3, 5, 7, 11, \ldots\}$
is a primitive subset of $\mathbb{N}_{\geq 2}$.
\end{lemma}
% [PROVE IT: Easy]
% SKETCH: Two distinct primes p != q cannot satisfy p | q, since the only
%   positive divisors of a prime q are 1 and q, and p > 1, p != q.

\begin{lemma}[Subsets of primitive sets are primitive]
\label{lem:subset-primitive}
If $A$ is primitive and $B \subseteq A$, then $B$ is primitive.
\end{lemma}
% [PROVE IT: Easy]
% SKETCH: Any pair in B is also a pair in A, so the non-divisibility
%   property is inherited.

\section{The Prime Case (Mertens)}

\begin{lemma}[Mertens' second theorem]
\label{lem:mertens}
As $x \to \infty$,
\[
  \sum_{p \leq x,\, p \text{ prime}} \frac{1}{p} \;=\; \log \log x + M + O\!\left(\frac{1}{\log x}\right),
\]
where $M \approx 0.2615$ is the Meissel--Mertens constant.
\end{lemma}
% [PROVE IT: Hard]
% SKETCH: Standard PNT-adjacent estimate. Use partial summation against
%   theta(x) = sum_{p <= x} log p, together with the PNT-style bound
%   theta(x) = x + O(x / log x), to convert sum 1/p into an integral
%   form and evaluate.

\begin{proposition}[Primes satisfy the bound]
\label{prop:primes-bound}
For primes $\mathbb{P}_x := \mathbb{P} \cap [x, \infty)$,
\[
  S(\mathbb{P}_x) \;=\; \sum_{p \geq x} \frac{1}{p \log p} \;=\; O\!\left(\frac{1}{\log x}\right)
  \qquad \text{as } x \to \infty.
\]
In particular $S(\mathbb{P}_x) = o(1)$.
\end{proposition}
% [PROVE IT: Medium]
% SKETCH: Partial summation. Let pi(t) be the prime counting function;
%   use pi(t) = O(t / log t) (Chebyshev/PNT) and integrate 1/(t log^2 t)
%   from x to infinity to get a 1/log x bound.

\section{Erdős' Boundedness Theorem (1935)}

\begin{theorem}[Erdős, 1935]
\label{thm:erdos-1935}
There exists an absolute constant $C > 0$ such that for every primitive
set $A \subseteq \mathbb{N}_{\geq 2}$,
\[
  S(A) \;=\; \sum_{a \in A} \frac{1}{a \log a} \;\leq\; C.
\]
\end{theorem}
% [PROVE IT: Hard]
% SKETCH: Classical argument. For each a in A, associate the "multiplicative
%   interval" M_a = { n in N : a | n, n/a has all prime factors < p(a) },
%   where p(a) is the smallest prime factor of a. The M_a are pairwise
%   disjoint (this uses primitivity of A), and sum_{n in M_a} 1/n is
%   comparable to 1/(a log a) via Mertens applied to the primes < p(a).
%   Summing over a and bounding by the total harmonic sum up to N gives
%   S(A) = O(1).

\begin{corollary}[Finite primitive sum]
\label{cor:finite-sum}
For every primitive set $A \subseteq \mathbb{N}_{\geq 2}$ the sum
$\sum_{a \in A} 1/(a \log a)$ converges.
\end{corollary}
% [PROVE IT: Easy]

\section{The Erdős--Sárközy--Szemerédi Conjecture}

\begin{conjecture}[Erdős--Sárközy--Szemerédi; now Theorem, Price 2026]
\label{thm:erdos-1196}
For every $x \geq 2$ and every primitive set $A \subseteq [x, \infty)$,
\[
  S(A) \;=\; \sum_{a \in A} \frac{1}{a \log a} \;\leq\; 1 + O\!\left(\frac{1}{\log x}\right).
\]
Consequently, $\sup_{A \subset [x,\infty) \text{ primitive}} S(A) = 1 + o(1)$
as $x \to \infty$.
\end{conjecture}
% [PROVE IT: Hard]
% SKETCH: This is the main Erdős Problem #1196. The idea (per Price's
%   recent argument) is to reduce to a "critical" family of primitive sets
%   parametrised by a Markov-style branching structure on the anatomy
%   of integers, and then estimate the supremum of S over that family
%   using Mertens (Lemma~\ref{lem:mertens}) and the Chebyshev bound on
%   pi(x). The prime case (Proposition~\ref{prop:primes-bound}) is the
%   extremal configuration: primes achieve the 1 + O(1/log x) rate and
%   any deviation from the prime-like structure decreases S.

\begin{corollary}[Matching lower bound]
\label{cor:lower-bound}
The bound in Theorem~\ref{thm:erdos-1196} is tight: the set of primes
in $[x, \infty)$ gives $S(\mathbb{P}_x) = 1 - o(1)$ is \emph{not} achieved,
but the family of "primes times a fixed integer" approaches $1$.
\end{corollary}
% [PROVE IT: Medium]
% SKETCH: For k >= 1 fixed, the set A_k = { kp : p prime, kp >= x }
%   is primitive (distinct primes do not divide each other, and the
%   common factor k is shared). Evaluate S(A_k) using Mertens and
%   compare against 1 as k varies. The supremum over such families
%   realises the constant 1.

\end{document}
